% Part: normal-modal-logic
% Chapter: completeness
% Section: lindenbaums-lemma

\documentclass[../../../include/open-logic-section]{subfiles}

\begin{document}

\olfileid{nml}{com}{lin}

\olsection{Lindenbaum 引理}

林登鲍姆引理断言，每个 $\Sigma$ 一致的公式集都包含在至少一个\emph{完全的} $\Sigma$ 一致集中。我们构造典范模型时将表明，对于每个完全的 $\Sigma$ 一致集~$\Delta$，在典范模型中都有一个世界，使得 $\Delta$ 中的且仅有 $\Delta$ 中的公式为真。因此，林登鲍姆引理保证了每个 $\Sigma$ 一致集在典范模型的某个世界中为真。

\begin{thm}[Lindenbaum 引理]\ollabel{thm:lindenbaum}
  如果 $\Gamma$ 是 $\Sigma$ 一致的，那么存在一个扩张~$\Gamma$ 的完全 $\Sigma$ 一致集 $\Delta$。
\end{thm}

\begin{proof}
令 $!A_0$, $!A_1$, \dots{} 为该语言所有公式的一个穷尽列举（允许重复）。例如，可以先列出 $\Obj p_0$，然后在每个阶段 $n \ge 1$，列出所有长度为~$n$、只使用变元 $\Obj p_0$, \dots,~$\Obj p_n$ 的有穷多个公式。我们通过对~$n$ 归纳来定义公式集 $\Delta_n$，然后令 $\Delta = \bigcup_n \Delta_n$。首先令 $\Delta_0 = \Gamma$。假设 $\Delta_n$ 已定义，我们按下式定义 $\Delta_{n+1}$：
\[
\Delta_{n+1} =
\begin{cases}
  \Delta_n \cup \{!A_n\}, & \text{如果 $\Delta_n \cup \{ !A_n\}$
    是 $\Sigma$ 一致的；} \\
  \Delta_n \cup \{ \lnot !A_n\}, & \text{否则。}
\end{cases}
\]
现在令 $\Delta = \bigcup_{n=0}^\infty \Delta_n$。

我们需要证明这个定义确实产生了一个具有所需性质的集合 $\Delta$，即 $\Gamma \subseteq \Delta$ 且 $\Delta$ 是完全 $\Sigma$ 一致的。

显然 $\Gamma \subseteq \Delta$，因为由构造有 $\Delta_0 \subseteq \Delta$，而 $\Delta_0 = \Gamma$。事实上，对所有~$n$ 都有 $\Delta_n \subseteq \Delta$，因为 $\Delta$ 是所有~$\Delta_n$ 的并集。（由于在构造的每一步，我们都向已构造的集合添加一个公式，$\Delta_n \subseteq
\Delta_{n+1}$，而 $\subseteq$ 是传递的，因此只要 $n \le m$ 就有 $\Delta_n \subseteq \Delta_{m}$。）在构造的每个阶段，我们要么加入 $!A_n$ 要么加入 $\lnot !A_n$，并且每个公式都（至少一次）出现在所有~$!A_n$ 的列表中。所以，对于每个 $!A$，要么 $!A \in \Delta$，要么 $\lnot !A \in \Delta$，因此根据定义 $\Delta$ 是完全的。

最后，我们需要证明 $\Delta$ 是 $\Sigma$ 一致的。为此，我们证明 如果 $\Delta$ 是 $\Sigma$ 不一致的，那么某个 $\Delta_n$ 将是 $\Sigma$ 不一致的，以及 所有 $\Delta_n$ 都是 $\Sigma$ 一致的。

假设 $\Delta$ 是 $\Sigma$ 不一致的。那么 $\Delta \Proves[\Sigma] \lfalse$，即存在 $!A_1$, \dots,~$!A_k \in \Delta$ 使得 $\Sigma \Proves !A_1 \lif (!A_2 \lif \cdots (!A_k
\lif \lfalse)\dots)$。由于 $\Delta = \bigcup_{n=0}^\infty \Delta_n$，每个 $!A_i \in \Delta_{n_i}$ 对某个~$n_i$ 成立。令 $n$ 为这些指标中最大的那个。因为 $n_i \le n$，$\Delta_{n_i} \subseteq \Delta_n$。所以，所有 $!A_i$ 都在某个~$\Delta_n$ 中。这将意味着 $\Delta_n \Proves[\Sigma] \lfalse$，即 $\Delta_n$ 是 $\Sigma$ 不一致的。

为证明每个 $\Delta_n$ 是 $\Sigma$ 一致的，我们对~$n$ 进行简单的归纳。$\Delta_0 = \Gamma$，而我们假设了 $\Gamma$ 是 $\Sigma$ 一致的。所以命题对 $n = 0$ 成立。现在假设命题对 $n$ 成立，即 $\Delta_n$ 是 $\Sigma$ 一致的。$\Delta_{n+1}$ 要么是 $\Delta_n \cup \{!A_n\}$（如果它是 $\Sigma$ 一致的），否则就是 $\Delta_n \cup \{\lnot!A_n\}$。在第一种情况下，$\Delta_{n+1}$ 显然是 $\Sigma$ 一致的。然而，根据 \olref[prf][con]{prop:consistencyfacts}\olref[prf][con]{prop:consistencyfacts-c}，要么 $\Delta_n \cup \{!A_n\}$ 要么 $\Delta_n \cup \{\lnot!A_n\}$ 是 $\Sigma$ 一致的，所以在另一种情况下 $\Delta_{n+1}$ 也是 $\Sigma$ 一致的。
\end{proof}

\begin{cor}\ollabel{cor:provability-characterization}
  $\Gamma \Proves[\Sigma] !A$ 当且仅当对于每个扩张 $\Gamma$ 的完全 $\Sigma$ 一致集 $\Delta$（包括 $\Gamma = \emptyset$ 的情形，此时我们得到模态系统 $\Sigma$ 的另一种刻画），都有 $!A \in \Delta$。
\end{cor}

\begin{proof}
  设 $\Gamma \Proves[\Sigma] !A$，并令 $\Delta$ 为任意一个扩张 $\Gamma$ 的完全 $\Sigma$ 一致集。
  若 $!A \notin \Delta$，则由极大性知 $\lnot!A \in \Delta$，
  于是 $\Delta \Proves[\Sigma] !A$（由单调性）且 $\Delta \Proves[\Sigma] \lnot!A$（由自反性），
  从而 $\Delta$ 是不一致的。
  反之，若 $\Gamma \Proves/[\Sigma] !A$，则 $\Gamma \cup \{ \lnot!A\}$ 是 $\Sigma$ 一致的，
  根据林登鲍姆引理，存在一个扩张 $\Gamma \cup \{ \lnot!A \}$ 的完全 $\Sigma$ 一致集 $\Delta$。
  由一致性，$!A \notin \Delta$。
\end{proof}

\end{document}
